%! TEX root = ukuran gambar.tex
% the main TeX file which is intended to compile, :VimtexReload after adjustment
   % this file should be included in the main file
% :h vimtex-tex-root
\documentclass[../menggambarteknik.tex]{subfiles}
%ensure the class of docs

\begin{document}

\section{Ukuran Gambar}

\subsection*{Tujuan}
Setelah mempelajari bahan dalam bab ini, seharusnya Anda dapat:
\begin{enumerate}
	\item Mengidentifikasi macam-macam cara pemberian ukuran.
	\item Meletakan pemberian ukuran yang benar pada gambar.
	\item Memodifikasi pemberian ukuran pada gambar teknik.
	\item Membedakan pemberian ukuran pada bagian dalam dan luar.
	\item Membaca dan menafsirkan ukuran-ukuran yang ada pada gambar kerja.
\end{enumerate}

\subsection{Macam-Macam Cara Pemberian Ukuran}
Untuk menunjukkan panjang, lebar, tinggi atau diameter benda, maka pada gambar dicantumkan ukurannya. Ukuran yang tercantum ini bisa yang sesungguhnya, tetapi jika benda yang digambar diperbesar atau diperkecil, maka dapat menggunakan skala. Pemberian ukuran pada gambar mesin tidak bisa dibuat sembarangan melainkan mengikuti aturan-aturan yang sudah ditetapkan.

Penyusunan ukuran pada gambar kerja, dapat dilakukan dengan beberapa macam cara ukuran, yaitu: ukuran berantai, ukuran sejajar, ukuran berurutan dan ukuran gabungan (lihat Gambar 23).

\begin{figure}
	\begin{center}
		\includegraphics[width=.95\textwidth]{pemberian ukuran.jpg}
	\end{center}
	\caption{Macam-macam cara pemberian ukuran}
	\label{fig:pemberianukuran}
\end{figure}


\subsection{Aturan Memberi Ukuran}
Penempatan angka ukuran pada gambar kerja mengikuti prosedur sebagai berikut: dilakukan 1 mm di atas garis ukur, ditengah-tengah dan teratur, angka dan garis ukuran harus terbaca baik horisontal maupun vertikal, ukuran-ukuran kecil (di bawah 10 mm) tanda panahnya ditempatkan di luar arah ukur dan ukurannya dicantumkan di atas atau disamping tanda panah ukuran, serta pengukuran dimulai dari basis yang terkecil hingga yang terbesar. Semua ukuran dalam gambar teknik mesin dalam satuan mm, dan tidak perlu dicantumkan satuannya, apabila ukuran dalam satuan yang lain, maka satuannya dicantumkan (misal inchi).

Untuk membatasi bagian yang diberi ukuran pada ujung garis ukurnya diberi anak panah. Perbandingan ukuran panjang dan lebar anak panah adalah 3 : 1 dan dihitamkan (lihat Gambar 24). Jika jarak antara dua garis lebih kecil dari 7 mm, garis ukuran pada kedua sisinya diperpanjang kemudian gambar panahnya diberikan sebelah luar, sedangkan untuk ukuran yang saling merapat dapat digunakan titik sebagai pengganti anak panah.

\begin{figure}
	\begin{center}
		\includegraphics[width=.95\textwidth]{tandapanah.png}
	\end{center}
	\caption{Tanda Anak Panah}
	\label{fig:tandaanakpanah}
\end{figure}

Untuk menulis ukuran-ukuran pada gambar kerja dilaukan sebagai berikut: gambarlah angka-angka ukuran dengan jelas, angka-angka ukuran digambarkan sedemikian, sehingga dapat dibaca dari sebelah bawah dan kanan dari gambar, dan ukuran ditempatkan sedemikian di mana bentuk atau profil dari potongan kerja diperlihatkan paling jelas.

Untuk angka ukuran yang tidak horisontal maupun vertikal, penempatannya diatur sedemikian sesuai dengan garis ukurnya. Ada daerah-daerah yang sebaiknya dihindari untuk penempatan angka ukuran, yaitu pada daerah 30o sebelah kiri bagian atas garis vertikal dan 30o bagian sebelah kanan garis vertikal bawah, pada Gambar 25 adalah daerah yang diarsir.

\begin{figure}
	\begin{center}
		\includegraphics[width=.95\textwidth]{angka bentuk miring.jpg}
	\end{center}
	\caption{Penempatan angka ukuran pada bentuk miring}
	\label{fig:angkamiring}
\end{figure}


Adanya ukuran pada gambar yang dibuat sesuai dengan aturan akan memperjelas bagi pembaca gambar tentang benda yang sebenarnya, tetapi ada kalanya ukuran yang berlebihan justru akan membingungkan. Untuk itu penunjukkan ukuran sebaiknya tidak berulang-ulang (hanya sekali).

Pada gambar, penunjukkan ukuran seluruhnya seharusnya diberikan agar mempermudah dalam menentukan kebutuhan bahan dari benda yang dibuat oleh pekerja. Ukuran seluruhnya (jumlah) bisa menjadi ukuran pembantu, tetapi bisa juga menjadi ukuran yang penting. Pada Gambar~\ref{fig:ukuranbantu} diperlihatkan di mana pada gambar a menunjukkan ukuran jumlah sebagai ukuran pembantu, sedangkan pada gambar b menunjukkan ukuran jumlah sebagai ukuran penting (pokok). Di dalam penunjukkannya ukuran pembantu ditulis di dalam kurung.


\begin{figure}
	\begin{center}
		\includegraphics[width=.95\textwidth]{ukuran bantu.png}
	\end{center}
	\caption{Ukuran seluruhnya sebagai ukuran bantu dan pokok}
	\label{fig:ukuranbantu}
\end{figure}


Pada gambar yang ditunjukkan bagian dalamnya, pemberian ukuran dipisahkan antara bagian luar dan bagian dalamnya. Untuk itu apabila ukuran bagian luar ditempatkan bagian atas, maka ukuran bagian dalam ditempatkan pada bagian bawah, demikian sebaliknya Gambar~\ref{fig:ukuranluardlm}.

\begin{figure}
	\begin{center}
		\includegraphics[width=.95\textwidth]{ukuranluardalam.jpg}
	\end{center}
	\caption{Penempatan ukuran bagian luar dan dalam}
	\label{fig:ukuranluardlm}
\end{figure}


% Ukuran untuk bagian yang berbentuk bulat, segi empat, dan sebagianya, apabila tidak dibuat gambar pandangan samping, atas atau bawah, maka pada pengukurannya perlu diberi lambang untuk bagian permukaan tersebut, dengan Ø dan $\square$. Gambar 28 menunjukkan penulisan ukuran lambang-lambang tersebut.

\begin{figure}
	\begin{center}
		
\includegraphics[width=.95\textwidth]{lambang bulan.png}
	\end{center}
	\caption{Ukuran dengan lambang bulat dan segi empat}
	\label{fig:lambangbulat}
\end{figure}


Angka-angka ukuran pada sumbu dan arsiran tidak boleh terpotong. Sumbu dan arsir dihilangkan pada angka ukuran yang dimaksud (lihat Gambar 29 di bawah ini).


\begin{figure}
	\begin{center}
		\includegraphics[width=.95\textwidth]{angka arsir.jpg}
	\end{center}
	\caption{Angka ukuran dalam arsir}
	\label{fig:angka arsir}
\end{figure}



Penulisan ukuran yang sama, bisa dibuat satu dengan mencantumkan jumlah, kemudian ditulis ukurannya. Misal ada lubang yang berdiameter 10 mm dan jumlahnya 8 buah, maka penulisan ukurannya dapat dilakukan seperti terlihat pada Gambar 30 di bawah ini.


\begin{figure}
	\begin{center}
		\includegraphics[width=.95\textwidth]{ukuran lubang.jpg}
	\end{center}
	\caption{Penulisan ukuran lubang dengan diameter sama}
	\label{fig:ukuran lubang}
\end{figure}




\bibliographystyle{apalike}
\bibliography{../filename.bib} % required for citecompletion, not work in mainfile
\end{document}

